% =====================================================================
%  2.1 Likninger og formler — fasit
% =====================================================================

\fk{Del 1 — Uttrykk, likning, identitet og funksjon}
\fo{1.1} a) uttrykk \ b) likning \ c) funksjon \ d) identitet
\fo{1.2} a) Ja: $3\cdot4-5=7$ \ b) Nei: $9 \neq 10$ \ c) Ja: $16 = 16$
\fo{1.3} a) Alle tall er løsninger \ b) Ingen løsning \ c) $x = 5$
\fo{1.4} a) Nei. For $x = 1$ er venstre side $36$ og høyre side $26$. Riktig er $x^2 + 10x + 25$.
b) Ja. Venstre side er $x^2 - 16$ etter konjugatsetningen.
c) Nei. For $x = 3$ er venstre side $11$ og høyre side $9$. Likningen er bare sann for $x = 1$ og $x = 2$.
\fo{1.5} a) $6x-2 = 6x-2$, så $0=0$: alle tall \ b) $3x-2 = 3x+4$, så $0 = 6$: ingen løsning \ c) $5x - 8 = 3x + 6$, så $x = 7$: én løsning
\fo{1.6} Alle linjene til og med $0 = 0$ er riktige. Feilen er siste pil: $0 = 0$ betyr at alle tall er løsninger, ikke at $x = 0$. Likningen er en identitet.
\fo{1.7} a) Alltid (for tillatte verdier): $(x^2+12x+36) - (x^2-12x+36) = 24x$. Identitet.\par
b) Noen ganger: $x(x-2) = 0$ gir $x = 0$ og $x = 2$.\par
c) Aldri: $x^2 \ge 0$, så $x^2 + 1 \ge 1$.
\fo{1.8} Eksempler, flere svar er riktige. a) $\frac{x}{2} + 1 = x - \frac{1}{2}$ \ b) $\frac{x}{3} + 1 = \frac{x}{3} + 2$ gir $0 = 1$ \ c) $\frac{x}{2} + \frac{x}{2} = x$ gir $0 = 0$
\fo{1.9} Nei. For $x = -1$ er venstre side $1$ og høyre side $-1$. Likningen er sann for alle $x \ge 0$, men ikke for negative $x$. En identitet må stemme for alle tillatte $x$. Likningen er altså noen ganger sann.
\fo{Sjekk 1} a) Uttrykk, det er ikke noe likhetstegn \ b) $4 = -1$ er usant: ingen løsning \ c) Vi kan bare prøve noen tall, og det er uendelig mange. Neste tall kan svikte. Algebra viser at det stemmer for alle.

\fk{Del 2 — Likninger med brøk og formler}
\fo{2.1} a) $x = 9$ \ b) $x = 20$ \ c) $x = 7$
\fo{2.2} a) $3(x+3) + 2x = 24$ \ b) $5x + 9 = 24$, så $x = 3$. Kontroll: $\frac{6}{2} + \frac{3}{3} = 4$.
\fo{2.3} a) $F$ og $C$ er variable. $1{,}8$ og $32$ er konstanter. \ b) $F = 77$ \ c) $C = 10$, fordi $18 = 1{,}8C$
\fo{2.4} a) $s = v \cdot t$ \ b) $t = \frac{s}{v}$
\fo{2.5} a) $2(2x-1) - (x+4) = 12$, $3x - 6 = 12$, $x = 6$\par
b) $5x - 2(x-3) = 30$, $3x + 6 = 30$, $x = 8$\par
c) $3(x+3) - 4(x-2) = 6$, $-x + 17 = 6$, $x = 11$
\fo{2.6} Feilen er tredje pil: $-2\cdot(-2) = +4$, ikke $-4$. Riktig: $3x + 18 - 2x + 4 = 24$, $x + 22 = 24$, $x = 2$. Kontroll: $4 - 0 = 4$.
\fo{2.7} a) $b = \frac{O}{2} - l$ \ b) $u = \frac{2s}{t} - v$ \ c) $h = \frac{3V}{G}$
\fo{2.8} a) $P = 15 + 2t$. $P$ og $t$ er variable, $15$ og $2$ er konstanter. \ b) 40 minutter \ c) $t = \frac{P-15}{2}$. Med $P = 95$: $t = \frac{80}{2} = 40$.
\fo{2.9} Marte har rett. Ganger vi med 6, får vi $2(x-4) = x-4$, så $x = 4$. Eli delte på $x-4$, og det er 0 når $x = 4$. Vi kan ikke dele på 0, og da mistet hun løsningen. Kontroll: begge sider er $0$.
\fo{2.10} a) Gang med $ab$: $bx + ax = ab$. Sett $x$ utenfor: $x(a+b) = ab$. Da er $x = \frac{ab}{a+b}$ ($a + b \neq 0$).\par
b) Er $b = -a$, blir venstre side $\frac{x}{a} - \frac{x}{a} = 0$, og vi får $0 = 1$. Ingen løsning.
\fo{Sjekk 2} a) $x + 4 = 2x$, så $x = 4$ \ b) $a = \frac{F}{m}$ \ c) $3(x+7) = 3x + 21$. Hele telleren skal ganges med 3, ikke bare $x$.

\fk{Ekstra utfordring}
\fo{U1} Gang med 6: $6x - 2(x-2) = 3(x+4) - 6$, så $4x + 4 = 3x + 6$ og $x = 2$. Kontroll: $2 - 0 = 2$ og $3 - 1 = 2$.
\fo{U2} Omskrevet er likningen $(a-2)x = a + 4$.
a) For $a = 3$ er $3x - 3 = 2x + 4$, så $x = 7$.\par
b) $a = 2$ gir $0 = 6$: ingen løsning.\par
c) Nei. Da måtte både $a - 2 = 0$ og $a + 4 = 0$, og det er umulig.
\fo{U3} $C = 1{,}8C + 32$ gir $-0{,}8C = 32$, så $C = -40$. Begge viser $-40$.
\fo{U4} a) $P_A = 49 + 16x$ og $P_B = 89 + 12x$.
b) $49 + 16x = 89 + 12x$ gir $4x = 40$, så $x = 10$ km. Begge koster 209 kr.\par
c) Grafene krysser i $(10, 209)$. 6 km: A koster 145 kr og B 161 kr, så velg A. 25 km: A koster 449 kr og B 389 kr, så velg B.
\fo{U5} a) Første hull: \lstinline[style=laereverk,numbers=none,frame=none]|b == d|. Andre hull: «ingen losning». Er $a = c$, forsvinner $x$. Da er $b = d$ en sann påstand (alle tall), og $b \neq d$ en usann påstand (ingen løsning).
b) \lstinline[style=laereverk,numbers=none,frame=none]|los(3, 1, 1, 5)| skriver «en losning: 2.0». \lstinline[style=laereverk,numbers=none,frame=none]|los(2, 1, 2, 4)| skriver «ingen losning».
\fo{U6} a) Fellesnevner $x(x+1)$: $\frac{(x+1) - x}{x(x+1)} = \frac{1}{x(x+1)}$ ($x \neq 0$ og $x \neq -1$).\par
b) Tre tall er bare tre av uendelig mange. Et bevis må gjelde for alle $x$, og det viser regningen i a).
\fo{U7} For $x \ge 0$ er $\sqrt{x^2} = x$, så likningen er sann for uendelig mange $x$. For $x = -1$ er venstre side $1$ og høyre side $-1$. Da er den ikke sann for alle $x$, og dermed ikke en identitet.
